\documentclass[11pt]{article}
\usepackage[utf8]{inputenc}
\usepackage{graphicx}
\usepackage{booktabs}
\usepackage{longtable}
\usepackage{hyperref}

\title{Supplementary Material\\
\large Evaluating Hallucinations in Large Language Model Responses to Proteomics Queries}

\author{}
\date{}

\begin{document}

\maketitle

\section*{Supplementary Methods}

\subsection*{S1. Query Generation Protocol}
Detailed methodology for generating the 1,000 proteomics queries used in this study, including stratification by complexity and domain.

\subsection*{S2. Expert Annotation Guidelines}
Complete annotation protocol provided to the three expert annotators, including training materials and adjudication procedures.

\subsection*{S3. Statistical Analysis Details}
Full specification of Bayesian Beta-Binomial models, prior distributions, and convergence diagnostics.

\section*{Supplementary Results}

\subsection*{S4. Complete Model Performance Metrics}
Extended tables showing all performance metrics for each model across all query categories.

\subsection*{S5. Domain-Specific Error Analysis}
Detailed breakdown of hallucination types by proteomics domain.

\subsection*{S6. Calibration Analysis}
Full calibration curves and ECE/MCE metrics for all models.

\section*{Supplementary Figures}

\begin{figure}[h]
\centering
\includegraphics[width=0.8\textwidth]{figures/output/supplementary_model_comparison.png}
\caption{\textbf{Supplementary Figure 1.} Extended model performance comparison.}
\end{figure}

\section*{Supplementary Tables}

\begin{table}[h]
\centering
\caption{\textbf{Supplementary Table 1.} Complete dataset characteristics.}
\begin{tabular}{lc}
\toprule
Characteristic & Value \\
\midrule
Total queries & 1,000 \\
Models evaluated & 5 \\
Expert annotators & 3 \\
Inter-rater agreement ($\kappa$) & 0.87 \\
\bottomrule
\end{tabular}
\end{table}

\end{document}
